\section{余下账本事件导航：eastern-jin-sixteen-kingdoms}

\label{sec:remaining-event-anchors-eastern-jin-sixteen-kingdoms}

本节为尚未在正文设置目标的账本记录建立直接跳转。锚点显示可核日期与事件名称；其解释仍应回到账本所列来源、置信提示及既有正文的区域和材料边界。

\subsection{事件导航与来源边界}

\eventanchor{JH-LEDGER-031-5}{太兴元年三月（318年）·“司马睿称帝，是为元帝，东晋以建康为都正式确立”的后续制度或军政调整}

\eventanchor{JH-LEDGER-031-6}{太兴元年三月（318年）·“司马睿称帝，是为元帝，东晋以建康为都正式确立”的异源材料与影响范围的复核}

\eventanchor{JH-LEDGER-032-1}{太兴二年（319年）·“石勒在襄国称赵王并建立独立政权，后赵由汉赵将领集团分出”的前期动员与资源准备}

\eventanchor{JH-LEDGER-032-2}{太兴二年（319年）·“石勒在襄国称赵王并建立独立政权，后赵由汉赵将领集团分出”的命令、联盟或地方响应}

\eventanchor{JH-LEDGER-032-3}{太兴二年（319年）·石勒在襄国称赵王并建立独立政权，后赵由汉赵将领集团分出}

\eventanchor{JH-LEDGER-032-4}{太兴二年（319年）·“石勒在襄国称赵王并建立独立政权，后赵由汉赵将领集团分出”的执行中的空间与人员处置}

\eventanchor{JH-LEDGER-032-5}{太兴二年（319年）·“石勒在襄国称赵王并建立独立政权，后赵由汉赵将领集团分出”的后续制度或军政调整}

\eventanchor{JH-LEDGER-032-6}{太兴二年（319年）·“石勒在襄国称赵王并建立独立政权，后赵由汉赵将领集团分出”的异源材料与影响范围的复核}

\eventanchor{JH-LEDGER-033-1}{光初二年（319年）·“刘曜在长安改汉国号为赵，以关中为中心重建政权身份”的前期动员与资源准备}

\eventanchor{JH-LEDGER-033-2}{光初二年（319年）·“刘曜在长安改汉国号为赵，以关中为中心重建政权身份”的命令、联盟或地方响应}

\eventanchor{JH-LEDGER-033-3}{光初二年（319年）·刘曜在长安改汉国号为赵，以关中为中心重建政权身份}

\eventanchor{JH-LEDGER-033-4}{光初二年（319年）·“刘曜在长安改汉国号为赵，以关中为中心重建政权身份”的执行中的空间与人员处置}

\eventanchor{JH-LEDGER-033-5}{光初二年（319年）·“刘曜在长安改汉国号为赵，以关中为中心重建政权身份”的后续制度或军政调整}

\eventanchor{JH-LEDGER-033-6}{光初二年（319年）·“刘曜在长安改汉国号为赵，以关中为中心重建政权身份”的异源材料与影响范围的复核}

\eventanchor{JH-LEDGER-034-1}{大兴三年（320年）·“杜弢领导的荆湘叛乱被王敦等镇压，东晋重新取得长江中游部分军政控制”的前期动员与资源准备}

\eventanchor{JH-LEDGER-034-2}{大兴三年（320年）·“杜弢领导的荆湘叛乱被王敦等镇压，东晋重新取得长江中游部分军政控制”的命令、联盟或地方响应}

\eventanchor{JH-LEDGER-034-3}{大兴三年（320年）·杜弢领导的荆湘叛乱被王敦等镇压，东晋重新取得长江中游部分军政控制}

\eventanchor{JH-LEDGER-034-4}{大兴三年（320年）·“杜弢领导的荆湘叛乱被王敦等镇压，东晋重新取得长江中游部分军政控制”的执行中的空间与人员处置}

\eventanchor{JH-LEDGER-034-5}{大兴三年（320年）·“杜弢领导的荆湘叛乱被王敦等镇压，东晋重新取得长江中游部分军政控制”的后续制度或军政调整}

\eventanchor{JH-LEDGER-034-6}{大兴三年（320年）·“杜弢领导的荆湘叛乱被王敦等镇压，东晋重新取得长江中游部分军政控制”的异源材料与影响范围的复核}
